\documentclass[14pt,usepdftitle=false,aspectratio=169]{beamer}
\usepackage{preambule}
\setbeamercolor{structure}{fg=black}
\usepackage{polynomes,al,usuelles,topologie}
\hypersetup{pdftitle={Théorie spectrale -- Théorie de Gelfand}}
\title{Théorie spectrale\\\emph{Théorie de Gelfand}}
\author{}
\date{}
\begin{document}
\begin{frame}
    \titlepage
\end{frame}
\slideq{Rayon spectral}{1/5}
\slider{$r_A\l a\r=\max[\lambda\in\sigma_A\l a\r]{\left|\lambda\right|}$}{1/5}
\slideq{Résolvante de $a$}{2/5}
\slider{$\appl{R_a}{\bbC\setminus\sigma\l a\r}{A^\times}{\lambda}{\l a-\lambda\r^{-1}}$\linebreak C'est une application holomorphe et $R_a'=R_a^2$}{2/5}
\slideq{Propriétés topologiques de $\sigma\l a\r$ pour $a\in A$ une algèbre de Banach}{3/5}
\slider{C'est un compact non vide contenu dans la boule de rayon $\left\|a\right\|$}{3/5}
\slideq{Théorème de Gelfand-Mazur}{4/5}
\slider{Soit $A$ une algèbre de Banach, alors pour tout $a\in A$, $r\l a\r=\lim[n\to+\infty]{\left\|a^n\right\|^\oldfrac1n}=\inf[n\to+\infty]{\left\|a^n\right\|^\oldfrac1n}$}{4/5}
\slideq{Propriétés topologiques de $A^\times$ pour $A$ une algèbre de Banach}{5/5}
\slider{$A^\times$ est un ouvert et est un groupe topologique}{5/5}
\end{document}